\documentclass[appendix_C]{subfiles}

\begin{document}
\section*{Приложение C. Реализация}
\addcontentsline{toc}{section}{Приложение C. Реализация}

\begin{lstlisting}
// Файл: ./src/basis/SimplexBasis.cpp

#include "basis/SimplexBasis.h"
#include <stdexcept>
#include <iostream>

SimplexBasis::SimplexBasis(const std::vector<Eigen::RowVectorXd>& vertices) : vertices_(vertices) {
    dim_ = (int)vertices.size() - 1;
    if ((int)vertices[0].size() != dim_)
        throw std::runtime_error("Несоответствие размерности вершин");

    Eigen::MatrixXd A(dim_ + 1, dim_ + 1);
    for (int i = 0; i <= dim_; ++i) {
        for (int j = 0; j < dim_; ++j) {
            A(i, j) = vertices[i][j];
        }
        A(i, dim_) = 1.0;
    }

    Ainv_ = A.inverse();

    std::cout << "A = \n" << A << "\n" << std::endl;
    std::cout << "Ainv = \n" << Ainv_ << "\n" << std::endl;
}

int SimplexBasis::size() const {
    return dim_ + 1;
}

double SimplexBasis::value(int i, const Eigen::RowVectorXd& x) const {
    if (x.size() != dim_)
        throw std::runtime_error("Несоответствие размерности точки");

    Eigen::RowVectorXd aug(dim_ + 1);
    for (int k = 0; k < dim_; ++k)
        aug(k) = x(k);
    aug(dim_) = 1.0;

    double res = 0.0;
    for (int k = 0; k <= dim_; ++k) {
        res += aug(k) * Ainv_(k, i);
    }
    return res;
}

bool SimplexBasis::verifyInterpolation(const SimplexBasis& basis, double tolerance) {
    const auto& vertices = basis.getVertices();
    double      maxError = 0.0;

    std::cout << "Проверка интерполяции: ℓ_i(v_j) должно быть 1 при i=j и 0 иначе" << std::endl;

    for (int i = 0; i < basis.size(); ++i) {
        for (int j = 0; j < basis.size(); ++j) {
            double val      = basis.value(i, vertices[j]);
            double expected = (i == j) ? 1.0 : 0.0;
            double error    = std::abs(val - expected);
            if (error > maxError)
                maxError = error;
            // Отладка: выводим первые несколько
            if (i < 3 && j < 3) {
                std::cout << "ℓ_" << i << "(v_" << j << ") = " << val;
                std::cout << " (ожидалось " << expected << ")" << std::endl;
            }
        }
    }

    std::cout << "Максимальная ошибка по всем парам (вершина, базис): " << maxError << std::endl;
    std::cout << std::endl;

    return maxError < tolerance;
}

=========================================================
// Файл: ./src/cli/CommandLineParser.cpp

#include "cli/CommandLineParser.h"
#include <iostream>
#include <string>

ProgramOptions CommandLineParser::parse(int argc, char* argv[]) {
    ProgramOptions opts;

    for (int i = 1; i < argc; ++i) {
        std::string arg = argv[i];

        if (arg == "-h" || arg == "--help") {
            opts.showHelp = true;
            return opts;
        }
        if (arg == "--dimension" && i + 1 < argc) {
            parseDimension(opts, i, argv);
            opts.endVertex = 1LL << opts.dimension;
        }
        if (arg == "--range" && i + 2 < argc) {
            parseRange(opts, i, argv);
        }
        if (arg == "--threads" && i + 1 < argc) {
            parseThreads(opts, i, argv);
        }
        if (arg == "--log-interval" && i + 1 < argc) {
            parseLogInterval(opts, i, argv);
        }
        if (arg == "--batch" && i + 1 < argc) {
            opts.batchMode = true;
            opts.inputDir  = argv[++i];
            std::cout << "Пакетный режим. Каталог: " << opts.inputDir << std::endl;
        }
        if (arg == "--load-hadamard" && i + 1 < argc) {
            opts.loadFromFile = true;
            opts.fileType     = ProgramOptions::FileType::HADAMARD;
            opts.inputFile    = argv[++i];
        }
    }

    return opts;
}

void CommandLineParser::parseRange(ProgramOptions& opts, int& i, char* argv[]) {
    opts.partial     = true;
    opts.startVertex = std::stoll(argv[++i]);
    opts.endVertex   = std::stoll(argv[++i]);
    std::cout << "Режим: частичный перебор" << std::endl;
    std::cout << "Диапазон вершин: [" << opts.startVertex << ", " << opts.endVertex << ")"
              << std::endl;
}

void CommandLineParser::parseDimension(ProgramOptions& opts, int& i, char* argv[]) {
    opts.dimension = std::stoi(argv[++i]);
    opts.endVertex = 1LL << opts.dimension;
    std::cout << "Размерность установлена: " << opts.dimension << std::endl;
    std::cout << "Всего вершин: " << opts.endVertex << " (2^" << opts.dimension << ")" << std::endl;
}

void CommandLineParser::parseThreads(ProgramOptions& opts, int& i, char* argv[]) {
    opts.numThreads = std::stoi(argv[++i]);
    std::cout << "Количество потоков: " << opts.numThreads << std::endl;
}

void CommandLineParser::parseLogInterval(ProgramOptions& opts, int& i, char* argv[]) {
    opts.logInterval = std::stoll(argv[++i]);
    std::cout << "Интервал логирования: " << opts.logInterval << " вершин" << std::endl;
}

void CommandLineParser::printHelp(const char* programName) {
    std::cout << "Использование: " << programName << " [ОПЦИИ]" << std::endl;
    std::cout << std::endl;
    std::cout << "ОПЦИИ РАЗМЕРНОСТИ:" << std::endl;
    std::cout << "  --dimension N       установить размерность пространства (по умолчанию 7)"
              << std::endl;
    std::cout << std::endl;
    std::cout << "ОПЦИИ ПЕРЕБОРА:" << std::endl;
    std::cout << "  --range START END   частичный перебор вершин в диапазоне [START, END)"
              << std::endl;
    std::cout << std::endl;
    std::cout << "ОПЦИИ ВЫПОЛНЕНИЯ:" << std::endl;
    std::cout << "  --threads N         количество потоков (0 = автоматически)" << std::endl;
    std::cout
        << "  --log-interval N    интервал логирования (количество вершин, по умолчанию 100000)"
        << std::endl;
    std::cout << std::endl;
    std::cout << "ОПЦИИ ЗАГРУЗКИ ДАННЫХ:" << std::endl;
    std::cout << "  --load-hadamard FILE  загрузить матрицу Адамара из файла" << std::endl;
    std::cout << "  --batch DIR           пакетная обработка всех файлов в каталоге" << std::endl;
    std::cout << std::endl;
    std::cout << "ПРИМЕРЫ:" << std::endl;
    std::cout << "  " << programName
              << "                                              # стандартный режим (dim=7)"
              << std::endl;
    std::cout << "  " << programName << " --dimension 31                               # dim=31"
              << std::endl;
    std::cout << "  " << programName << " --dimension 31 --threads 16                  # 16 потоков"
              << std::endl;
    std::cout << "  " << programName
              << " --dimension 31 --log-interval 1000000        # логирование каждые 1 млн вершин"
              << std::endl;
    std::cout << "  " << programName
              << " --range 0 500000000                          # частичный перебор" << std::endl;
    std::cout << "  " << programName
              << " --load-hadamard ../preset/test_Hadamar_8.txt # загрузка матрицы Адамара"
              << std::endl;
}

=========================================================
// Файл: ./src/hadamard_utils/HadamardMatrixIterator.cpp

#include "hadamard_utils/HadamardMatrixIterator.h"
#include <sstream>
#include <stdexcept>
#include <iostream>
#include <cctype>
#include <cmath>

HadamardMatrixIterator::HadamardMatrixIterator(const std::string& filename)
    : filename_(filename), currentOrder_(0) {
    reset();
}

void HadamardMatrixIterator::reset() {
    file_.close();
    file_.open(filename_);
    if (!file_.is_open()) {
        throw std::runtime_error("Не удалось открыть файл: " + filename_);
    }
}

bool HadamardMatrixIterator::hasNext() const {
    return file_.is_open() && !file_.eof();
}

bool HadamardMatrixIterator::isCommentLine(const std::string& line) {
    if (line.empty())
        return false;
    char c = line[0];
    // Строка комментария: первый символ не +, -, 0, 1, и не минус (который может быть началом -1)
    // Но -1 начинается с '-', поэтому нужно отличать числовой -1 от комментария
    // Проще: комментарий — если первый символ буква или не является частью допустимого формата
    if (c == '+' || c == '-')
        return false;
    if (c == '0' || c == '1')
        return false;
    if (c == ' ' || c == '\t')
        return false;
    return true;
}

bool HadamardMatrixIterator::isSymbolicLine(const std::string& line) {
    if (line.empty())
        return false;
    // Символьный формат: строка состоит только из '+' и '-'
    for (char ch : line) {
        if (ch != '+' && ch != '-')
            return false;
    }
    return true;
}

bool HadamardMatrixIterator::isNumericLine(const std::string& line) {
    if (line.empty())
        return false;
    // Числовой формат: строка содержит числа -1, 0 или 1, разделённые пробелами
    std::istringstream iss(line);
    double             val;
    while (iss >> val) {
        if (std::abs(val) > 1.01)
            return false;  // не -1, 0 или 1
    }
    return true;
}

std::vector<double> HadamardMatrixIterator::parseLine(const std::string& line) {
    std::vector<double> row;

    if (isSymbolicLine(line)) {
        for (char ch : line) {
            if (ch == '+')
                row.push_back(1.0);
            else if (ch == '-')
                row.push_back(-1.0);
        }
    } else {
        std::istringstream iss(line);
        double             val;
        while (iss >> val) {
            // Нормализуем к -1, 0, 1 (на случай -0.999999)
            if (std::abs(val - 1.0) < 1e-12)
                val = 1.0;
            else if (std::abs(val + 1.0) < 1e-12)
                val = -1.0;
            else if (std::abs(val) < 1e-12)
                val = 0.0;
            row.push_back(val);
        }
    }

    return row;
}

Eigen::MatrixXd HadamardMatrixIterator::readNextMatrix() {
    std::vector<std::vector<double>> rows;

    std::string line;
    while (std::getline(file_, line)) {
        // Удаляем пробелы
        size_t start = line.find_first_not_of(" \t\r\n");
        if (start == std::string::npos) {
            // Пустая строка — разделитель матриц, выходим
            break;
        }
        line = line.substr(start);

        size_t end = line.find_last_not_of(" \t\r\n");
        if (end != std::string::npos) {
            line = line.substr(0, end + 1);
        }

        if (line.empty())
            break;

        if (isCommentLine(line)) {
            continue;
        }

        // Строка данных матрицы
        if (isSymbolicLine(line) || isNumericLine(line)) {
            std::vector<double> row = parseLine(line);
            rows.push_back(row);
        } else {
            // Неизвестный формат — вероятно, разделитель
            break;
        }
    }

    int n = rows.size();
    if (n == 0) {
        throw std::runtime_error("Не удалось прочитать матрицу");
    }

    int m = rows[0].size();
    if (n != m) {
        throw std::runtime_error("Матрица не квадратная: " + std::to_string(n) + "x" +
                                 std::to_string(m));
    }

    Eigen::MatrixXd H(n, n);
    for (int i = 0; i < n; ++i) {
        for (int j = 0; j < n; ++j) {
            H(i, j) = rows[i][j];
        }
    }

    currentOrder_ = n;
    return H;
}

Eigen::MatrixXd HadamardMatrixIterator::next() {
    return readNextMatrix();
}

=========================================================
// Файл: ./src/hadamard_utils/HadamardUtils.cpp

#include "hadamard_utils/HadamardUtils.h"

#include <iostream>

Eigen::MatrixXd HadamardUtils::hadamardTo01Matrix(const Eigen::MatrixXd& H) {
    int             n = H.rows();
    Eigen::MatrixXd B = H;

    // Шаг 1: каждый столбец, начинающийся с -1, умножается на -1
    for (int j = 0; j < n; ++j) {
        if (B(0, j) == -1.0) {
            for (int i = 0; i < n; ++i) {
                B(i, j) = -B(i, j);
            }
        }
    }

    // Шаг 2: каждая строка, начинающаяся с -1, умножается на -1
    for (int i = 0; i < n; ++i) {
        if (B(i, 0) == -1.0) {
            for (int j = 0; j < n; ++j) {
                B(i, j) = -B(i, j);
            }
        }
    }

    // Шаг 3: подматрица C (строки и столбцы 2..n), замена: 1 → 0, -1 → 1
    int             dim = n - 1;
    Eigen::MatrixXd D(dim, dim);

    for (int i = 0; i < dim; ++i) {
        for (int j = 0; j < dim; ++j) {
            double val = B(i + 1, j + 1);
            D(i, j)    = (val == -1.0) ? 1.0 : 0.0;
        }
    }

    return D;
}

std::vector<Eigen::RowVectorXd> HadamardUtils::buildVerticesFrom01Matrix(const Eigen::MatrixXd& D) {
    int dim         = D.rows();
    int numVertices = dim + 1;

    std::vector<Eigen::RowVectorXd> vertices;
    vertices.reserve(numVertices);

    // Нулевая вершина
    Eigen::RowVectorXd zero(dim);
    zero.setZero();
    vertices.push_back(zero);

    // Строки матрицы D как вершины
    for (int i = 0; i < dim; ++i) {
        vertices.push_back(D.row(i));
    }

    return vertices;
}

std::vector<Eigen::RowVectorXd> HadamardUtils::getVerticesWithLogging(const Eigen::MatrixXd& H,
                                                                      int                    dim) {
    int order = dim + 1;

    std::cout << "=== Матрица Адамара H (1/-1): ===\n" << H << "\n" << std::endl;

    long double detH           = H.determinant();
    long double theoreticalDet = std::pow(order, order / 2);
    std::cout << "det(H) = " << detH << std::endl;
    std::cout << "Теоретическое значение: ±" << theoreticalDet << "\n" << std::endl;

    // Проверка определителя матрицы Адамара
    if (std::abs(std::abs(detH) - std::abs(theoreticalDet)) > 10) {
        std::cout << "ОШИБКА: Определитель матрицы Адамара (-1/1) не соответствует теоретическому "
                     "значению!"
                  << std::endl;
        std::cout << "Получено: " << detH << ", ожидается ±" << theoreticalDet << std::endl;
        std::exit(1);
    }

    Eigen::MatrixXd D = hadamardTo01Matrix(H);

    std::cout << "=== Матрица D (0/1): ===\n" << D << "\n" << std::endl;

    double detD        = D.determinant();
    double expectedDet = std::abs(detH) / std::pow(2.0, dim);
    std::cout << "det(D) = " << detD << std::endl;
    std::cout << "Теоретическое значение: ±" << expectedDet << "\n" << std::endl;

    // Проверка определителя (0/1)-матрицы
    if (std::abs(std::abs(detD) - std::abs(expectedDet)) > 10) {
        std::cout << "ОШИБКА: Определитель матрицы Адамара (0/1) не соответствует теоретическому "
                     "значению!"
                  << std::endl;
        std::cout << "Получено: " << detD << ", ожидается ±" << expectedDet << std::endl;
        std::exit(1);
    }

    auto vertices = buildVerticesFrom01Matrix(D);

    std::cout << "Вершины симплекса:" << std::endl;
    for (size_t i = 0; i < vertices.size(); ++i) {
        std::cout << "v[" << i << "]: " << vertices[i] << std::endl;
    }

    std::cout << std::endl;

    return vertices;
}

bool HadamardUtils::isHadamardMatrix(const Eigen::MatrixXd& H, double tolerance) {
    int n = H.rows();

    // Шаг 1: проверка квадратности
    if (H.cols() != n) {
        std::cerr << "Ошибка: матрица не квадратная" << std::endl;
        return false;
    }

    // Шаг 2: проверка элементов (±1)
    for (int i = 0; i < n; ++i) {
        for (int j = 0; j < n; ++j) {
            double val = std::abs(H(i, j));
            if (std::abs(val - 1.0) > tolerance) {
                std::cerr << "Ошибка: элемент [" << i << "][" << j << "] = " << H(i, j)
                          << " не равен ±1" << std::endl;
                return false;
            }
        }
    }

    // Шаг 3: проверка ортогональности строк (H * H^T = n * I)
    Eigen::MatrixXd HHT = H * H.transpose();
    for (int i = 0; i < n; ++i) {
        for (int j = 0; j < n; ++j) {
            double expected = (i == j) ? n : 0.0;
            if (std::abs(HHT(i, j) - expected) > tolerance) {
                std::cerr << "Ошибка: (H*H^T)[" << i << "][" << j << "] = " << HHT(i, j)
                          << ", ожидается " << expected << std::endl;
                return false;
            }
        }
    }

    return true;
}
=========================================================
// Файл: ./src/processing/Processing.cpp

#include "processing/Processing.h"

#include "basis/SimplexBasis.h"
#include "cli/CommandLineParser.h"
#include "vertices_provider/SimplexProviderFactory.h"
#include "upper_bound_calculator/FullEnumUpperBoundCalculator.h"
#include "upper_bound_calculator/LebesgueFunction.h"
#include "hadamard_utils/HadamardMatrixIterator.h"
#include "hadamard_utils/HadamardUtils.h"

#include <chrono>
#include <iomanip>
#include <iostream>
#include <filesystem>

namespace fs = std::filesystem;

struct ProgramOptions;

// Функция обработки одной матрицы Адамара
double Processing::processHadamardMatrix(const Eigen::MatrixXd& H, const ProgramOptions& opts,
                                         int matrixCounter) {
    int order         = H.rows();
    int expectedOrder = opts.dimension + 1;

    // Проверка порядка матрицы
    if (order != expectedOrder) {
        std::cout << "Предупреждение: порядок матрицы " << order << " не соответствует ожидаемому "
                  << expectedOrder << std::endl;
        std::cout << "Матрица #" << matrixCounter << ": неверный порядок (" << order << " вместо "
                  << expectedOrder << ")" << std::endl;
        return -1.0;
    }

    std::cout << "\n" << std::string(60, '=') << std::endl;
    std::cout << "Обработка матрицы #" << matrixCounter << " (порядок " << order << ")"
              << std::endl;
    std::cout << std::string(60, '=') << std::endl;

    // Строим симплекс из матрицы Адамара
    auto vertices = HadamardUtils::getVerticesWithLogging(H, order - 1);

    SimplexBasis     basis(vertices);
    LebesgueFunction lebesgueFunc(basis);

    // Проверка интерполяции
    if (!SimplexBasis::verifyInterpolation(basis, 0.00001)) {
        std::cout << "ОШИБКА: Интерполяционные свойства не выполняются!" << std::endl;
        std::cout << "Матрица #" << matrixCounter << ": ошибка интерполяции" << std::endl;
        return -1.0;
    }

    // Вычисление верхней границы
    FullEnumUpperBoundCalculator calculator(lebesgueFunc, opts.dimension);
    calculator.setNumThreads(opts.numThreads);
    calculator.setLogInterval(opts.logInterval);

    auto   start       = std::chrono::high_resolution_clock::now();
    double maxLebesgue = calculator.calculate();
    auto   end         = std::chrono::high_resolution_clock::now();
    auto   duration    = std::chrono::duration_cast<std::chrono::seconds>(end - start);

    Eigen::RowVectorXd maxPoint = calculator.getMaxPoint();

    std::cout << "Максимальная сумма модулей: " << std::fixed << std::setprecision(6)
              << maxLebesgue << std::endl;
    std::cout << "Время выполнения: " << duration.count() << " секунд" << std::endl;
    std::cout << std::endl;

    // Вывод результата
    std::cout << "ИТОГ для Матрица #" << matrixCounter << ", порядок " << order << ": "
              << std::fixed << std::setprecision(6) << maxLebesgue << " (за " << duration.count()
              << " сек)" << std::endl;

    return maxLebesgue;
}

// Функция обработки одного файла (может содержать несколько матриц)
int Processing::processHadamardMatrixFromFile(const std::string&    filename,
                                              const ProgramOptions& opts, double& globalMin,
                                              int&                globalMinMatrixIdx,
                                              Eigen::RowVectorXd& globalMinPoint) {
    int                matrixCounter = 1;
    double             fileMin       = std::numeric_limits<double>::max();
    int                fileMinIdx    = -1;
    Eigen::RowVectorXd fileMinPoint;

    std::cout << "\n" << std::string(80, '=') << std::endl;
    std::cout << "Обработка файла: " << filename << std::endl;
    std::cout << std::string(80, '=') << std::endl;

    try {
        HadamardMatrixIterator it(filename);
        while (it.hasNext()) {
            Eigen::MatrixXd H        = it.next();
            double          lebesgue = processHadamardMatrix(H, opts, matrixCounter);

            if (lebesgue >= 0) {
                if (lebesgue < fileMin) {
                    fileMin    = lebesgue;
                    fileMinIdx = matrixCounter;
                }
            }
            matrixCounter++;
        }
    } catch (const std::exception& e) {
        std::cout << "Ошибка при обработке файла " << filename << ": " << e.what() << std::endl;
        std::cout << "Ошибка в файле " << filename << ": " << e.what() << std::endl;
    }

    if (fileMinIdx != -1) {
        std::cout << "\n>>> МИНИМУМ в файле " << filename << ": матрица #" << fileMinIdx
                  << " дала значение " << std::fixed << std::setprecision(4) << fileMin
                  << std::endl;

        // Обновляем глобальный минимум
        if (fileMin < globalMin) {
            globalMin          = fileMin;
            globalMinMatrixIdx = fileMinIdx;
            globalMinPoint     = fileMinPoint;
        }
    }

    return 0;
}

int Processing::processBatch(const ProgramOptions& opts) {
    // Собираем все .txt файлы в каталоге
    std::vector<std::string> files;
    if (fs::exists(opts.inputDir) && fs::is_directory(opts.inputDir)) {
        for (const auto& entry : fs::directory_iterator(opts.inputDir)) {
            if (entry.is_regular_file() && entry.path().extension() == ".txt") {
                files.push_back(entry.path().string());
            }
        }
    } else {
        std::cerr << "Каталог не найден: " << opts.inputDir << std::endl;
        return 1;
    }

    std::sort(files.begin(), files.end());
    std::cout << "Найдено файлов для обработки: " << files.size() << std::endl;

    // Глобальный минимум
    double             globalMin          = std::numeric_limits<double>::max();
    int                globalMinMatrixIdx = -1;
    Eigen::RowVectorXd globalMinPoint;

    for (const auto& file : files) {
        processHadamardMatrixFromFile(file, opts, globalMin, globalMinMatrixIdx, globalMinPoint);
    }

    // Вывод глобального минимума
    if (globalMinMatrixIdx != -1) {
        std::cout << "\n" << std::string(80, '*') << std::endl;
        std::cout << "ГЛОБАЛЬНЫЙ МИНИМУМ: матрица #" << globalMinMatrixIdx << " дала значение "
                  << std::fixed << std::setprecision(6) << globalMin << std::endl;
        std::cout << std::string(80, '*') << std::endl;
    }

    return 0;
}

int Processing::processSingleFile(const std::string& filename, const ProgramOptions& opts) {
    double             dummyMin;
    int                dummyIdx;
    Eigen::RowVectorXd dummyPoint;
    return processHadamardMatrixFromFile(filename, opts, dummyMin, dummyIdx, dummyPoint);
}

int Processing::processSilvester(const ProgramOptions& opts) {
    auto vertices_provider =
        SimplexProviderFactory::create(SimplexProviderFactory::Type::SILVESTER, opts.dimension, "");
    auto vertices = vertices_provider->getVertices();

    std::cout << "Число вершин симплекса: " << vertices.size() << "\n" << std::endl;

    SimplexBasis     basis(vertices);
    LebesgueFunction lebesgueFunc(basis);

    std::cout << "\n=== Проверка интерполяции ===" << std::endl;
    if (!SimplexBasis::verifyInterpolation(basis, 0.00001)) {
        std::cout << "ВНИМАНИЕ: Интерполяционные свойства не выполняются!" << std::endl;
        return 1;
    }
    std::cout << "ВЕРНО: Интерполяционные свойства выполняются" << "\n" << std::endl;

    FullEnumUpperBoundCalculator calculator(lebesgueFunc, opts.dimension);
    calculator.setNumThreads(opts.numThreads);
    calculator.setLogInterval(opts.logInterval);

    auto   start       = std::chrono::high_resolution_clock::now();
    double maxLebesgue = calculator.calculate();
    auto   end         = std::chrono::high_resolution_clock::now();
    auto   duration    = std::chrono::duration_cast<std::chrono::seconds>(end - start);

    std::cout << "\n=== РЕЗУЛЬТАТЫ ===" << std::endl;
    std::cout << std::fixed << std::setprecision(6);
    std::cout << "Максимальная сумма модулей: " << maxLebesgue << std::endl;

    Eigen::RowVectorXd maxPoint = calculator.getMaxPoint();
    std::cout << "Найдена в вершине: [";
    for (int i = 0; i < maxPoint.size(); ++i) {
        std::cout << static_cast<int>(maxPoint(i) + 0.5);
        if (i + 1 < maxPoint.size())
            std::cout << " ";
    }
    std::cout << "]" << std::endl;

    std::cout << "Время выполнения: " << duration.count() << " секунд" << std::endl;

    std::cout << "Сильвестр (dim=" << opts.dimension << "): " << std::fixed << std::setprecision(6)
              << maxLebesgue << " (за " << duration.count() << " сек)" << std::endl;
    return 0;
}


=========================================================
// Файл: ./src/upper_bound_calculator/FullEnumUpperBoundCalculator.cpp

#include "upper_bound_calculator/FullEnumUpperBoundCalculator.h"
#include <thread>
#include <vector>
#include <iostream>
#include <iomanip>
#include <chrono>
#include <mutex>

static std::mutex cout_mutex;

FullEnumUpperBoundCalculator::FullEnumUpperBoundCalculator(const LebesgueFunction& func, int dim)
    : func_(func),
      dim_(dim),
      globalMax_(0.0),
      processedVertices_(0),
      lastLogged_(0),
      currentBestVertex_(0),
      startVertex_(0),
      endVertex_(0),
      partial_(false) {
    maxPoint_ = Eigen::RowVectorXd::Zero(dim_);
}

FullEnumUpperBoundCalculator::FullEnumUpperBoundCalculator(const LebesgueFunction& func, int dim,
                                                           int64_t startVertex, int64_t endVertex)
    : func_(func),
      dim_(dim),
      globalMax_(0.0),
      processedVertices_(0),
      lastLogged_(0),
      currentBestVertex_(0),
      startVertex_(startVertex),
      endVertex_(endVertex),
      partial_(true) {
    maxPoint_ = Eigen::RowVectorXd::Zero(dim_);
}

static void enumerateSubset(const LebesgueFunction* func, int dim, int64_t start, int64_t end,
                            std::atomic<double>* globalMax, std::atomic<int64_t>* processed,
                            std::atomic<int64_t>* lastLogged, std::atomic<int64_t>* bestVertex,
                            int64_t totalVertices, int64_t logInterval,
                            std::chrono::steady_clock::time_point startTime) {
    double  localMax        = 0.0;
    int64_t localBestVertex = 0;

    for (int64_t idx = start; idx < end; ++idx) {
        Eigen::RowVectorXd x(dim);
        for (int i = 0; i < dim; ++i) {
            x(i) = ((idx >> i) & 1) ? 1.0 : 0.0;
        }

        double val = (*func)(x);

        // Обновляем локальный и глобальный максимум сразу при нахождении лучшей вершины
        if (val > localMax) {
            localMax        = val;
            localBestVertex = idx;

            double currentMax = globalMax->load();
            while (localMax > currentMax) {
                if (globalMax->compare_exchange_weak(currentMax, localMax)) {
                    bestVertex->store(localBestVertex);
                    break;
                }
                currentMax = globalMax->load();  // Перечитываем при неудаче из-за другого потока
            }
        }

        int64_t processedCount = processed->fetch_add(1) + 1;

        // Логирование каждые logInterval итераций
        if (processedCount - lastLogged->load() >= logInterval) {
            lastLogged->store(processedCount);

            double percent =
                100.0 * static_cast<double>(processedCount) / static_cast<double>(totalVertices);
            auto now = std::chrono::steady_clock::now();
            auto elapsed =
                std::chrono::duration_cast<std::chrono::seconds>(now - startTime).count();
            auto elapsed_us =
                std::chrono::duration_cast<std::chrono::microseconds>(now - startTime).count();

            double us_per_iter =
                (processedCount > 0) ? static_cast<double>(elapsed_us) / processedCount : 0.0;
            double vertices_per_sec =
                (elapsed > 0) ? static_cast<double>(processedCount) / elapsed : 0.0;

            double remaining_sec = 0.0;
            if (processedCount > 0) {
                remaining_sec = static_cast<double>(elapsed) * (totalVertices - processedCount) /
                                processedCount;
            }

            auto time_t = std::chrono::system_clock::to_time_t(std::chrono::system_clock::now());

            std::lock_guard<std::mutex> lock(cout_mutex);
            std::cout << std::put_time(std::localtime(&time_t), "%H:%M:%S")
                      << " | Прогресс: " << std::fixed << std::setprecision(4) << percent << "%"
                      << " | Обработано: " << processedCount << " / " << totalVertices
                      << " | Прошло: " << elapsed << " сек"
                      << " | Скорость: " << std::setprecision(0) << vertices_per_sec << " в/сек"
                      << " | мкс/итер: " << std::setprecision(2) << us_per_iter
                      << " | Осталось ≈: " << static_cast<int>(remaining_sec) << " сек"
                      << " | Текущее: " << val << " | Рекорд: " << globalMax->load() << std::endl;
        }
    }
}

double FullEnumUpperBoundCalculator::calculate() {
    int64_t totalVertices = partial_ ? (endVertex_ - startVertex_) : (1LL << dim_);
    int64_t startOffset   = partial_ ? startVertex_ : 0;

    int numThreads = numThreads_;
    if (numThreads <= 0) {
        numThreads = std::thread::hardware_concurrency();
        if (numThreads == 0)
            numThreads = 8;
    }

    if (!partial_) {
        std::cout << "=== Полный перебор вершин гиперкуба ===" << std::endl;
        std::cout << "Размерность: " << dim_ << std::endl;
        std::cout << "Всего вершин: " << totalVertices << " (2^" << dim_ << ")" << std::endl;
    } else {
        std::cout << "=== Частичный перебор вершин гиперкуба ===" << std::endl;
        std::cout << "Размерность: " << dim_ << std::endl;
        std::cout << "Диапазон: [" << startVertex_ << ", " << endVertex_ << ")" << std::endl;
        std::cout << "Всего вершин: " << totalVertices << std::endl;
    }
    std::cout << "Потоков: " << numThreads << std::endl;
    std::cout << "Логирование каждые " << logInterval_ << " вершин" << std::endl;
    std::cout << "=========================================" << std::endl;

    auto startTime = std::chrono::steady_clock::now();

    int64_t                  chunkSize = totalVertices / numThreads;
    std::vector<std::thread> threads;

    for (int t = 0; t < numThreads; ++t) {
        int64_t start = startOffset + t * chunkSize;
        int64_t end   = (t == numThreads - 1) ? startOffset + totalVertices : start + chunkSize;

        threads.emplace_back(enumerateSubset, &func_, dim_, start, end, &globalMax_,
                             &processedVertices_, &lastLogged_, &currentBestVertex_, totalVertices,
                             logInterval_, startTime);
    }

    for (auto& th : threads) {
        th.join();
    }

    auto endTime  = std::chrono::steady_clock::now();
    auto duration = std::chrono::duration_cast<std::chrono::seconds>(endTime - startTime);

    std::cout << "\n=== Перебор завершён ===" << std::endl;
    std::cout << "Время: " << duration.count() << " секунд" << std::endl;
    std::cout << "Обработано вершин: " << processedVertices_.load() << std::endl;

    double avgSpeed = 0.0;
    if (duration.count() > 0) {
        avgSpeed = static_cast<double>(processedVertices_.load()) / duration.count();
    }
    std::cout << "Средняя скорость: " << std::fixed << std::setprecision(0) << avgSpeed << " в/сек"
              << std::endl;

    std::cout << std::fixed << std::setprecision(3);
    std::cout << "Максимальная сумма модулей: " << globalMax_.load() << std::endl;

    int64_t bestIdx = currentBestVertex_.load();
    for (int i = 0; i < dim_; ++i) {
        maxPoint_(i) = ((bestIdx >> i) & 1) ? 1.0 : 0.0;
    }

    return globalMax_.load();
}


=========================================================
// Файл: ./src/upper_bound_calculator/LebesgueFunction.cpp

#include "upper_bound_calculator/LebesgueFunction.h"
#include <cmath>

LebesgueFunction::LebesgueFunction(const SimplexBasis& basis) : basis_(basis) {
}

double LebesgueFunction::operator()(const Eigen::RowVectorXd& x) const {
    double sum = 0.0;
    for (int i = 0; i < basis_.size(); ++i) {
        sum += std::abs(basis_.value(i, x));
    }
    return sum;
}

=========================================================
// Файл: ./src/vertices_provider/HadamardFileProvider.cpp

#include "vertices_provider/HadamardFileProvider.h"

#include "hadamard_utils/HadamardMatrixIterator.h"
#include "hadamard_utils/HadamardUtils.h"
#include <fstream>
#include <iostream>
#include <sstream>
#include <stdexcept>

HadamardFileProvider::HadamardFileProvider(const std::string& filename)
    : filename_(filename), dim_(0) {
    H_ = readHadamardFile(filename_);

    int order = H_.rows();
    dim_      = order - 1;

    // Проверка, что загруженная матрица является матрицей Адамара
    if (!HadamardUtils::isHadamardMatrix(H_)) {
        throw std::runtime_error("Загруженная матрица не является матрицей Адамара!");
    }

    std::cout << "Матрица Адамара порядка " << order << " успешно загружена и проверена"
              << std::endl;
}

std::vector<Eigen::RowVectorXd> HadamardFileProvider::getVertices() {
    return HadamardUtils::getVerticesWithLogging(H_, dim_);
}

Eigen::MatrixXd HadamardFileProvider::readHadamardFile(const std::string& filename) {
    std::ifstream file(filename);
    if (!file.is_open()) {
        throw std::runtime_error("Не удалось открыть файл: " + filename);
    }

    std::vector<std::vector<double>> rows;
    std::string                      line;
    bool                             isSymbolicFormat = false;

    while (std::getline(file, line)) {
        // Удаляем пробелы в начале и конце
        line.erase(0, line.find_first_not_of(" \t\r\n"));
        line.erase(line.find_last_not_of(" \t\r\n") + 1);

        // Пропускаем пустые строки
        if (line.empty())
            continue;

        // Определяем формат по первому непустому символу
        char firstChar = line[0];

        if (firstChar == '+' || firstChar == '-') {
            isSymbolicFormat = true;

            // Парсим символьный формат
            std::vector<double> row;
            for (char ch : line) {
                if (ch == '+') {
                    row.push_back(1.0);
                } else if (ch == '-') {
                    row.push_back(-1.0);
                }
                // Игнорируем пробелы и другие символы
            }
            if (!row.empty()) {
                rows.push_back(row);
            }
        } else if (firstChar == '1' || (firstChar >= '0' && firstChar <= '9')) {
            isSymbolicFormat = false;

            std::istringstream  iss(line);
            std::vector<double> row;
            double              val;
            while (iss >> val) {
                row.push_back(val);
            }
            if (!row.empty()) {
                rows.push_back(row);
            }
        } else {
            // Строка-комментарий (начинается с буквы или другого символа)
            std::cout << "Пропущена строка-комментарий: " << line << std::endl;
        }
    }

    if (rows.empty()) {
        throw std::runtime_error("Файл не содержит данных");
    }

    int n = rows.size();
    int m = rows[0].size();

    if (n != m) {
        throw std::runtime_error("Матрица не квадратная: " + std::to_string(n) + "x" +
                                 std::to_string(m));
    }

    Eigen::MatrixXd H(n, n);
    for (int i = 0; i < n; ++i) {
        for (int j = 0; j < n; ++j) {
            H(i, j) = rows[i][j];
        }
    }

    std::cout << "Файл прочитан в " << (isSymbolicFormat ? "символьном" : "числовом") << " формате"
              << std::endl;

    return H;
}

HadamardMatrixIterator HadamardFileProvider::getIterator(const std::string& filename) {
    return HadamardMatrixIterator(filename);
}

=========================================================
// Файл: ./src/vertices_provider/SilvesterMethodProvider.cpp

#include "vertices_provider/SilvesterMethodProvider.h"
#include "hadamard_utils/HadamardUtils.h"
#include <iostream>
#include <stdexcept>
#include <cmath>

SilvesterMethodProvider::SilvesterMethodProvider(int dim) : dim_(dim) {
    int order = dim + 1;
    H_        = hadamardMatrix(order);

    if (!HadamardUtils::isHadamardMatrix(H_)) {
        throw std::runtime_error("Созданная алгоритмом матрица не является матрицей Адамара!");
    }
}

Eigen::MatrixXd SilvesterMethodProvider::hadamardMatrix(int order) {
    if (order == 1) {
        Eigen::MatrixXd H(1, 1);
        H(0, 0) = 1.0;
        return H;
    }
    if (order <= 0 || (order & (order - 1)) != 0) {
        throw std::runtime_error("Метод Сильвестра требует порядок, являющийся степенью двойки");
    }
    int             half  = order / 2;
    Eigen::MatrixXd Hhalf = hadamardMatrix(half);
    Eigen::MatrixXd H(order, order);
    H.topLeftCorner(half, half)     = Hhalf;
    H.topRightCorner(half, half)    = Hhalf;
    H.bottomLeftCorner(half, half)  = Hhalf;
    H.bottomRightCorner(half, half) = -Hhalf;
    return H;
}

std::vector<Eigen::RowVectorXd> SilvesterMethodProvider::getVertices() {
    return HadamardUtils::getVerticesWithLogging(H_, dim_);
}

=========================================================
// Файл: ./src/vertices_provider/SimplexProviderFactory.cpp

#include "vertices_provider/SimplexProviderFactory.h"
#include "vertices_provider/SilvesterMethodProvider.h"
#include "vertices_provider/HadamardFileProvider.h"

std::unique_ptr<ISimplexProvider> SimplexProviderFactory::create(Type type, int dim,
                                                                 const std::string& filename) {
    switch (type) {
        case Type::SILVESTER:
            return std::make_unique<SilvesterMethodProvider>(dim);
        case Type::FROM_HADAMARD_FILE:
            return std::make_unique<HadamardFileProvider>(filename);
        default:
            throw std::runtime_error("Неизвестный тип провайдера");
    }
}

\end{lstlisting}
\end{document}